\documentclass[11pt,a4paper]{article}
\usepackage[utf8]{inputenc}
\usepackage[margin=0.8in]{geometry}
\usepackage{listings}
\usepackage{xcolor}
\usepackage{titlesec}
\usepackage{hyperref}
\usepackage{fancyhdr}

\definecolor{codebg}{RGB}{245,245,245}
\definecolor{codegreen}{RGB}{0,128,0}
\definecolor{codegray}{RGB}{128,128,128}
\definecolor{codeblue}{RGB}{0,0,180}
\definecolor{codepurple}{RGB}{128,0,128}

\lstset{
    language=C++,
    backgroundcolor=\color{codebg},
    basicstyle=\ttfamily\footnotesize,
    keywordstyle=\color{codeblue}\bfseries,
    commentstyle=\color{codegreen}\itshape,
    stringstyle=\color{codepurple},
    numberstyle=\tiny\color{codegray},
    numbers=left,
    numbersep=5pt,
    frame=single,
    breaklines=true,
    breakatwhitespace=false,
    tabsize=4,
    showstringspaces=false,
    captionpos=b,
    morekeywords={override,nullptr,string,shared_ptr,unique_ptr,make_shared,make_unique,
                  lock_guard,unique_lock,mutex,thread,vector,unordered_map,map,
                  chrono,steady_clock,time_point,atomic},
}

\pagestyle{fancy}
\fancyhf{}
\fancyhead[L]{\textbf{LLD --- Meeting Scheduler}}
\fancyhead[R]{\thepage}

\title{\Huge\textbf{Meeting Scheduler}\\[0.3em]\Large Low-Level Design --- C++ Concurrency}
\author{LLD Interview Preparation}
\date{}

\begin{document}
\maketitle
\tableofcontents
\newpage

% ============================================================
\section{File Structure}
% ============================================================
\begin{verbatim}
meeting-scheduler/
  models.hpp       -- Interval, Participant (participantMtx),
                      MeetingRoom (roomMtx), Meeting
  strategies.hpp   -- RoomSelectionStrategy, MeetingObserver
  managers.hpp     -- MeetingScheduler (Singleton)
  main.cpp         -- Demo: sequential, concurrent, cancellation,
                      strategy switch
\end{verbatim}

\textbf{Compilation:}
\begin{verbatim}
g++ -std=c++17 -pthread -o meeting_scheduler.exe main.cpp
\end{verbatim}

\subsection{Concurrency Summary}
\begin{itemize}
    \item \textbf{Level 1:} \texttt{schedulerMtx} (MeetingScheduler) --- meetings, rooms, participants
    \item \textbf{Level 2:} \texttt{roomMtx} (per-MeetingRoom) --- room bookings
    \item \textbf{Level 3:} \texttt{participantMtx} (per-Participant) --- participant schedule
    \item \textbf{Lock order:} always L1 $\rightarrow$ L2 $\rightarrow$ L3 (never reverse)
    \item Release L1 before per-resource locks where possible
    \item Rollback on participant conflict (undo booked participants)
\end{itemize}

\subsection{Design Patterns}
\begin{itemize}
    \item \textbf{Singleton:} MeetingScheduler
    \item \textbf{Strategy:} BestFitRoomStrategy, FirstFitRoomStrategy, LargestRoomStrategy
    \item \textbf{Observer:} EmailNotifier, CalendarSyncObserver
\end{itemize}

\subsection{Smart Pointers}
\begin{itemize}
    \item \texttt{unique\_ptr<MeetingRoom>} --- scheduler sole owner
    \item \texttt{unique\_ptr<Participant>} --- scheduler sole owner
    \item \texttt{shared\_ptr<Meeting>} --- scheduler map + caller
    \item \texttt{unique\_ptr<RoomSelectionStrategy>} --- scheduler sole owner
    \item \texttt{shared\_ptr<MeetingObserver>} --- scheduler vector + caller
\end{itemize}

% ============================================================
\newpage
\section{models.hpp}
% ============================================================
\begin{lstlisting}[title={\texttt{models.hpp}}]
#ifndef MODELS_HPP
#define MODELS_HPP

#include <string>
#include <vector>
#include <iostream>
#include <mutex>
#include <memory>
#include <algorithm>
using namespace std;

// ============================================================
// INTERVAL
// ============================================================
class Interval {
    int start, end;
public:
    Interval(int s, int e) : start(s), end(e) {
        if (s >= e) throw invalid_argument("start must be < end");
    }
    int getStart() const { return start; }
    int getEnd() const { return end; }
    bool overlaps(const Interval& o) const {
        return start < o.end && o.start < end;
    }
};

// ============================================================
// PARTICIPANT
// ============================================================
class Participant {
    string id;
    string name;
    // Per-participant calendar with its own mutex
    // PESSIMISTIC: lock before checking/modifying schedule
    vector<Interval> schedule;
    mutable mutex participantMtx;

public:
    Participant(const string& id, const string& name)
        : id(id), name(name) {}

    string getId() const { return id; }
    string getName() const { return name; }

    // Check if participant is free during interval
    bool isFree(const Interval& interval) const {
        lock_guard<mutex> lock(participantMtx);
        for (const auto& s : schedule) {
            if (s.overlaps(interval)) return false;
        }
        return true;
    }

    // Atomically check and book slot -- returns false if conflict
    bool tryBook(const Interval& interval) {
        lock_guard<mutex> lock(participantMtx);
        for (const auto& s : schedule) {
            if (s.overlaps(interval)) return false;
        }
        schedule.push_back(interval);
        return true;
    }

    // Remove a booked slot (for cancellation)
    void removeSlot(const Interval& interval) {
        lock_guard<mutex> lock(participantMtx);
        schedule.erase(
            remove_if(schedule.begin(), schedule.end(),
                [&](const Interval& s) {
                    return s.getStart() == interval.getStart() &&
                           s.getEnd() == interval.getEnd();
                }),
            schedule.end());
    }

    void displaySchedule() const {
        lock_guard<mutex> lock(participantMtx);
        cout << "    " << name << ": ";
        if (schedule.empty()) { cout << "free" << endl; return; }
        for (const auto& s : schedule)
            cout << "[" << s.getStart() << "-" << s.getEnd() << "] ";
        cout << endl;
    }
};

// ============================================================
// MEETING ROOM
// Has its own mutex for per-room locking.
// Multiple rooms can be booked in parallel (no contention).
// ============================================================
class MeetingRoom {
    string id;
    int capacity;
    vector<Interval> bookings;
    mutable mutex roomMtx;

public:
    MeetingRoom(const string& id, int cap)
        : id(id), capacity(cap) {}

    string getId() const { return id; }
    int getCapacity() const { return capacity; }

    bool isAvailable(const Interval& interval) const {
        lock_guard<mutex> lock(roomMtx);
        for (const auto& b : bookings) {
            if (b.overlaps(interval)) return false;
        }
        return true;
    }

    // Atomic check-and-book
    bool tryBook(const Interval& interval) {
        lock_guard<mutex> lock(roomMtx);
        for (const auto& b : bookings) {
            if (b.overlaps(interval)) return false;
        }
        bookings.push_back(interval);
        return true;
    }

    void cancelSlot(const Interval& interval) {
        lock_guard<mutex> lock(roomMtx);
        bookings.erase(
            remove_if(bookings.begin(), bookings.end(),
                [&](const Interval& b) {
                    return b.getStart() == interval.getStart() &&
                           b.getEnd() == interval.getEnd();
                }),
            bookings.end());
    }

    void display() const {
        lock_guard<mutex> lock(roomMtx);
        cout << "  Room[" << id << "] cap=" << capacity << " bookings: ";
        if (bookings.empty()) { cout << "none" << endl; return; }
        for (const auto& b : bookings)
            cout << "[" << b.getStart() << "-" << b.getEnd() << "] ";
        cout << endl;
    }
};

// ============================================================
// MEETING
// Stores meeting info. No mutex needed -- immutable after creation,
// only modified under MeetingScheduler's lock for cancellation.
// ============================================================
enum class MeetingStatus { SCHEDULED, CANCELLED };

class Meeting {
    string meetingId;
    string organizerId;
    Interval interval;
    vector<string> participantIds;
    string roomId;
    MeetingStatus status;

public:
    Meeting(const string& id, const string& organizer, const Interval& interval,
            const vector<string>& participants, const string& room)
        : meetingId(id), organizerId(organizer), interval(interval),
          participantIds(participants), roomId(room),
          status(MeetingStatus::SCHEDULED) {}

    string getId() const { return meetingId; }
    string getOrganizerId() const { return organizerId; }
    const Interval& getInterval() const { return interval; }
    const vector<string>& getParticipantIds() const { return participantIds; }
    string getRoomId() const { return roomId; }
    MeetingStatus getStatus() const { return status; }
    void setStatus(MeetingStatus s) { status = s; }

    void display() const {
        cout << "  Meeting[" << meetingId << "] "
             << "[" << interval.getStart() << "-" << interval.getEnd() << "] "
             << "Room: " << roomId << " Organizer: " << organizerId
             << " Participants: ";
        for (const auto& p : participantIds) cout << p << " ";
        cout << "| " << (status == MeetingStatus::SCHEDULED ? "SCHEDULED" : "CANCELLED")
             << endl;
    }
};

#endif
\end{lstlisting}

% ============================================================
\newpage
\section{strategies.hpp}
% ============================================================
\begin{lstlisting}[title={\texttt{strategies.hpp}}]
#ifndef STRATEGIES_HPP
#define STRATEGIES_HPP

#include "models.hpp"
#include <vector>
#include <memory>
using namespace std;

// ============================================================
// STRATEGY PATTERN: Room Selection
// Different policies for how to pick a room.
// unique_ptr<RoomSelectionStrategy>: MeetingScheduler is sole owner.
// ============================================================

class RoomSelectionStrategy {
public:
    virtual ~RoomSelectionStrategy() = default;
    virtual MeetingRoom* selectRoom(vector<unique_ptr<MeetingRoom>>& rooms,
                                     int requiredCapacity,
                                     const Interval& interval) = 0;
    virtual string getName() const = 0;
};

// Pick smallest room that fits (minimize waste)
class BestFitRoomStrategy : public RoomSelectionStrategy {
public:
    MeetingRoom* selectRoom(vector<unique_ptr<MeetingRoom>>& rooms,
                             int requiredCapacity,
                             const Interval& interval) override {
        MeetingRoom* best = nullptr;
        int bestCap = INT_MAX;
        for (auto& r : rooms) {
            if (r->getCapacity() >= requiredCapacity &&
                r->getCapacity() < bestCap &&
                r->isAvailable(interval)) {
                best = r.get();
                bestCap = r->getCapacity();
            }
        }
        return best;
    }
    string getName() const override { return "BestFit"; }
};

// Pick first available room that fits
class FirstFitRoomStrategy : public RoomSelectionStrategy {
public:
    MeetingRoom* selectRoom(vector<unique_ptr<MeetingRoom>>& rooms,
                             int requiredCapacity,
                             const Interval& interval) override {
        for (auto& r : rooms) {
            if (r->getCapacity() >= requiredCapacity &&
                r->isAvailable(interval)) {
                return r.get();
            }
        }
        return nullptr;
    }
    string getName() const override { return "FirstFit"; }
};

// Pick largest room available (luxury mode)
class LargestRoomStrategy : public RoomSelectionStrategy {
public:
    MeetingRoom* selectRoom(vector<unique_ptr<MeetingRoom>>& rooms,
                             int requiredCapacity,
                             const Interval& interval) override {
        MeetingRoom* best = nullptr;
        int bestCap = 0;
        for (auto& r : rooms) {
            if (r->getCapacity() >= requiredCapacity &&
                r->getCapacity() > bestCap &&
                r->isAvailable(interval)) {
                best = r.get();
                bestCap = r->getCapacity();
            }
        }
        return best;
    }
    string getName() const override { return "Largest"; }
};

// ============================================================
// OBSERVER PATTERN: Meeting Notifications
// shared_ptr<MeetingObserver>: shared between scheduler's vector
// and caller (to remove later). Shared ownership.
// ============================================================

class MeetingObserver {
public:
    virtual ~MeetingObserver() = default;
    virtual void onMeetingScheduled(const Meeting& m) = 0;
    virtual void onMeetingCancelled(const Meeting& m) = 0;
    virtual string getName() const = 0;
};

class EmailNotifier : public MeetingObserver {
public:
    void onMeetingScheduled(const Meeting& m) override {
        cout << "    [Email] Meeting " << m.getId() << " scheduled "
             << "[" << m.getInterval().getStart() << "-"
             << m.getInterval().getEnd() << "] Room: " << m.getRoomId() << endl;
    }
    void onMeetingCancelled(const Meeting& m) override {
        cout << "    [Email] Meeting " << m.getId() << " cancelled" << endl;
    }
    string getName() const override { return "Email"; }
};

class CalendarSyncObserver : public MeetingObserver {
public:
    void onMeetingScheduled(const Meeting& m) override {
        cout << "    [Calendar] Synced meeting " << m.getId() << endl;
    }
    void onMeetingCancelled(const Meeting& m) override {
        cout << "    [Calendar] Removed meeting " << m.getId() << endl;
    }
    string getName() const override { return "CalendarSync"; }
};

#endif
\end{lstlisting}

% ============================================================
\newpage
\section{managers.hpp}
% ============================================================
\begin{lstlisting}[title={\texttt{managers.hpp}}]
#ifndef MANAGERS_HPP
#define MANAGERS_HPP

#include "strategies.hpp"
#include <unordered_map>
#include <mutex>
#include <thread>
using namespace std;

// ============================================================
// LOCKING: 3-level pessimistic
// Level 1: schedulerMtx (MeetingScheduler) -> meetings map, rooms, participants
// Level 2: roomMtx (per-MeetingRoom) -> room bookings
// Level 3: participantMtx (per-Participant) -> participant schedule
//
// LOCK ORDER: always L1 -> L2 -> L3 (never reverse)
// Release L1 before per-resource locks where possible.
//
// SMART POINTERS:
//   unique_ptr<MeetingRoom>     -> scheduler is sole owner of rooms
//   unique_ptr<Participant>     -> scheduler is sole owner of participants
//   shared_ptr<Meeting>         -> scheduler map + caller both hold ref
//   unique_ptr<RoomStrategy>    -> scheduler sole owner, swapped via move()
//   shared_ptr<MeetingObserver> -> scheduler vector + caller (to remove)
// ============================================================

class MeetingScheduler {
    static mutex singletonMtx;
    mutable mutex schedulerMtx;

    vector<unique_ptr<MeetingRoom>> rooms;
    unordered_map<string, unique_ptr<Participant>> participants;
    unordered_map<string, shared_ptr<Meeting>> meetings;
    unique_ptr<RoomSelectionStrategy> roomStrategy;
    vector<shared_ptr<MeetingObserver>> observers;
    int meetingCounter = 0;

    MeetingScheduler() {
        roomStrategy = make_unique<BestFitRoomStrategy>();
    }

    // Snapshot pattern: copy observers under lock, notify outside
    void notifyScheduled(const Meeting& m) {
        vector<shared_ptr<MeetingObserver>> snap;
        { lock_guard<mutex> lock(schedulerMtx); snap = observers; }
        for (auto& obs : snap) obs->onMeetingScheduled(m);
    }

    void notifyCancelled(const Meeting& m) {
        vector<shared_ptr<MeetingObserver>> snap;
        { lock_guard<mutex> lock(schedulerMtx); snap = observers; }
        for (auto& obs : snap) obs->onMeetingCancelled(m);
    }

public:
    static MeetingScheduler* getInstance() {
        lock_guard<mutex> lock(singletonMtx);
        static MeetingScheduler instance;
        return &instance;
    }

    void setRoomStrategy(unique_ptr<RoomSelectionStrategy> strategy) {
        lock_guard<mutex> lock(schedulerMtx);
        roomStrategy = move(strategy);
        cout << "[Scheduler] Room strategy: " << roomStrategy->getName() << endl;
    }

    void addRoom(const string& id, int capacity) {
        lock_guard<mutex> lock(schedulerMtx);
        rooms.push_back(make_unique<MeetingRoom>(id, capacity));
        cout << "[Scheduler] Room added: " << id << " (cap " << capacity << ")" << endl;
    }

    void addParticipant(const string& id, const string& name) {
        lock_guard<mutex> lock(schedulerMtx);
        participants[id] = make_unique<Participant>(id, name);
        cout << "[Scheduler] Participant added: " << name << endl;
    }

    void addObserver(shared_ptr<MeetingObserver> obs) {
        lock_guard<mutex> lock(schedulerMtx);
        observers.push_back(obs);
    }

    // ============================================================
    // SCHEDULE MEETING -- core concurrent operation
    // 1. Lock scheduler to generate ID, look up participants, select room
    // 2. Release scheduler lock
    // 3. Lock each participant (Level 3) to check/book their schedule
    // 4. Lock room (Level 2) to book the room
    // 5. Lock scheduler again to store meeting
    // ============================================================
    shared_ptr<Meeting> scheduleMeeting(const string& organizerId,
                                         const vector<string>& participantIds,
                                         int start, int end) {
        Interval interval(start, end);
        string meetingId;
        vector<Participant*> participantPtrs;
        MeetingRoom* room = nullptr;

        // Step 1: Under scheduler lock -- ID gen, lookups, room selection
        {
            lock_guard<mutex> lock(schedulerMtx);
            meetingId = "MTG_" + to_string(++meetingCounter);

            // Verify all participants exist
            for (const auto& pid : participantIds) {
                auto it = participants.find(pid);
                if (it == participants.end()) {
                    cout << "[Thread " << this_thread::get_id() << "] "
                         << "Participant not found: " << pid << endl;
                    return nullptr;
                }
                participantPtrs.push_back(it->second.get());
            }

            // Select room (strategy reads room availability -- room locks itself)
            int requiredCap = participantIds.size();
            room = roomStrategy->selectRoom(rooms, requiredCap, interval);
            if (!room) {
                cout << "[Thread " << this_thread::get_id() << "] "
                     << meetingId << " FAILED: no room available" << endl;
                return nullptr;
            }
        }
        // Scheduler lock released

        // Step 2: Check & book each participant (per-participant lock)
        vector<Participant*> booked;
        bool allFree = true;
        for (auto* p : participantPtrs) {
            if (p->tryBook(interval)) {
                booked.push_back(p);
            } else {
                allFree = false;
                cout << "[Thread " << this_thread::get_id() << "] "
                     << meetingId << " CONFLICT: " << p->getName()
                     << " busy [" << start << "-" << end << "]" << endl;
                break;
            }
        }

        // Rollback if any participant had a conflict
        if (!allFree) {
            for (auto* p : booked) p->removeSlot(interval);
            return nullptr;
        }

        // Step 3: Book the room (per-room lock)
        if (!room->tryBook(interval)) {
            // Room got taken between strategy check and booking (race)
            for (auto* p : booked) p->removeSlot(interval);
            cout << "[Thread " << this_thread::get_id() << "] "
                 << meetingId << " FAILED: room " << room->getId()
                 << " taken (race)" << endl;
            return nullptr;
        }

        // Step 4: Store meeting under scheduler lock
        auto meeting = make_shared<Meeting>(meetingId, organizerId, interval,
                                             participantIds, room->getId());
        {
            lock_guard<mutex> lock(schedulerMtx);
            meetings[meetingId] = meeting;
        }

        cout << "[Thread " << this_thread::get_id() << "] "
             << meetingId << " scheduled [" << start << "-" << end << "] "
             << "Room: " << room->getId() << endl;

        notifyScheduled(*meeting);
        return meeting;
    }

    // ============================================================
    // CANCEL MEETING
    // Release room slot, release participant slots, mark cancelled
    // ============================================================
    bool cancelMeeting(const string& meetingId) {
        shared_ptr<Meeting> meeting;
        {
            lock_guard<mutex> lock(schedulerMtx);
            auto it = meetings.find(meetingId);
            if (it == meetings.end()) {
                cout << "Meeting not found: " << meetingId << endl;
                return false;
            }
            meeting = it->second;
        }

        if (meeting->getStatus() == MeetingStatus::CANCELLED) {
            cout << "Meeting already cancelled: " << meetingId << endl;
            return false;
        }

        // Release room
        {
            lock_guard<mutex> lock(schedulerMtx);
            for (auto& r : rooms) {
                if (r->getId() == meeting->getRoomId()) {
                    r->cancelSlot(meeting->getInterval());
                    break;
                }
            }
        }

        // Release participants
        {
            lock_guard<mutex> lock(schedulerMtx);
            for (const auto& pid : meeting->getParticipantIds()) {
                auto it = participants.find(pid);
                if (it != participants.end()) {
                    it->second->removeSlot(meeting->getInterval());
                }
            }
        }

        meeting->setStatus(MeetingStatus::CANCELLED);
        cout << "[Thread " << this_thread::get_id() << "] "
             << meetingId << " cancelled" << endl;

        notifyCancelled(*meeting);
        return true;
    }

    void displayMeetings() {
        lock_guard<mutex> lock(schedulerMtx);
        cout << "\n=== All Meetings ===" << endl;
        for (auto& [id, m] : meetings) m->display();
    }

    void displayRooms() {
        lock_guard<mutex> lock(schedulerMtx);
        cout << "\n=== Rooms ===" << endl;
        for (auto& r : rooms) r->display();
    }

    void displayParticipants() {
        lock_guard<mutex> lock(schedulerMtx);
        cout << "\n=== Participants ===" << endl;
        for (auto& [id, p] : participants) p->displaySchedule();
    }

    MeetingScheduler(const MeetingScheduler&) = delete;
    MeetingScheduler& operator=(const MeetingScheduler&) = delete;
};

mutex MeetingScheduler::singletonMtx;

#endif
\end{lstlisting}

% ============================================================
\newpage
\section{main.cpp}
% ============================================================
\begin{lstlisting}[title={\texttt{main.cpp}}]
#include <iostream>
#include <thread>
#include <vector>
#include "managers.hpp"
using namespace std;

int main() {
    auto* sched = MeetingScheduler::getInstance();

    cout << "=== Meeting Scheduler ===" << endl << endl;

    // Observers
    sched->addObserver(make_shared<EmailNotifier>());
    sched->addObserver(make_shared<CalendarSyncObserver>());

    // Rooms (unique_ptr<MeetingRoom>: scheduler sole owner)
    sched->addRoom("R1", 5);
    sched->addRoom("R2", 10);
    sched->addRoom("R3", 3);

    // Participants (unique_ptr<Participant>: scheduler sole owner)
    sched->addParticipant("P1", "Alice");
    sched->addParticipant("P2", "Bob");
    sched->addParticipant("P3", "Charlie");
    sched->addParticipant("P4", "Diana");

    // ========== Sequential meetings ==========
    cout << "\n--- Sequential Meetings ---" << endl;
    sched->scheduleMeeting("P1", {"P1", "P2"}, 9, 10);         // R3 (best fit, cap 3 >= 2)
    sched->scheduleMeeting("P2", {"P2", "P3", "P4"}, 10, 12);  // R1 (best fit, cap 5 >= 3)
    sched->scheduleMeeting("P1", {"P1", "P2"}, 9, 11);         // FAIL: P1 & P2 busy 9-10

    // ========== Strategy switch ==========
    cout << "\n--- Switch to LargestRoom strategy ---" << endl;
    sched->setRoomStrategy(make_unique<LargestRoomStrategy>());
    sched->scheduleMeeting("P3", {"P3", "P4"}, 9, 10);  // Gets R2 (largest)

    // ========== Concurrent meetings ==========
    // 3 threads try to book overlapping slots
    // Per-participant locks detect conflicts atomically
    // Per-room locks prevent double-booking
    cout << "\n--- Concurrent Meetings (3 threads race for 14-15) ---" << endl;
    sched->setRoomStrategy(make_unique<FirstFitRoomStrategy>());

    vector<shared_ptr<Meeting>> results(3);
    vector<thread> threads;

    auto bookFn = [&](int idx, const string& org, vector<string> parts, int s, int e) {
        results[idx] = sched->scheduleMeeting(org, parts, s, e);
    };

    threads.emplace_back(bookFn, 0, "P1", vector<string>{"P1", "P2"}, 14, 15);
    threads.emplace_back(bookFn, 1, "P3", vector<string>{"P3", "P2"}, 14, 15); // P2 conflict
    threads.emplace_back(bookFn, 2, "P4", vector<string>{"P4"}, 14, 15);       // Different people, OK

    for (auto& t : threads) t.join();

    // ========== Cancellation ==========
    cout << "\n--- Cancellation ---" << endl;
    sched->cancelMeeting("MTG_1");
    // Now P1 and P2 are free 9-10, R3 is free 9-10
    cout << "\n--- Rebooking after cancellation ---" << endl;
    sched->scheduleMeeting("P3", {"P1", "P2"}, 9, 10);  // Should succeed now

    // ========== Display ==========
    sched->displayMeetings();
    sched->displayRooms();
    sched->displayParticipants();

    cout << "\n=== Done ===" << endl;
    return 0;
}
\end{lstlisting}

\end{document}
